\begin{lema}[Rango de las igualdades de línea]\label{lem:rango-llegada}
Un tipo realizable de \(v\) sitios distintos prescribe sus cuatro particiones en fibras de fila,
columna, diagonal y antidiagonal. Sean \(c_s\) sus componentes y \(D\) la dimensión del espacio
de posiciones. Si la componente \(g\) contiene \(v_g\) sitios y tiene número ciclomático
\(\beta_g\), su rango es \(v_g-1+L_g\), con \(L_g\ge\min(\beta_g,2)\).
En particular \(D=v+c_s-\sum_gL_g\le p-1\), para los tres exponentes de la
Proposición~\ref{prop:soportes-llegada}.
\end{lema}
\begin{proof}
Sustituimos cada fibra por un árbol sobre sus sitios. Sus aristas imponen igualdades
\(n_f\cdot(z_i-z_l)=0\), donde \(n_f\) es una de las cuatro normales. Se conservan las
igualdades, las componentes y \(E-V+c_s\). Un árbol generador de una componente da
\(v_g-1\) ecuaciones independientes por eliminación de hojas. Una dependencia en un
subgrafo asignaría coeficientes no nulos a sus aristas de modo que en cada vértice la
suma con signo de las normales ponderadas fuera cero. En su soporte no hay hojas.
En un vértice de grado dos, ambas aristas son de la misma recta. Suprimimos cada camino
de vértices de grado dos. Una componente hecha sólo de tales vértices, o un lazo tras la
supresión, sería un ciclo dentro del árbol de una recta. Dos aristas paralelas darían dos
caminos rectos entre los mismos sitios distintos, que determinan una única recta. También
cerrarían un ciclo en aquel árbol. El grafo reducido es simple y de grado mínimo tres.
Por componente, \(V\ge4\) y \(E-V+1\ge V/2+1\ge3\). Suprimir los caminos conserva el
número de ciclos. Por tanto un árbol generador con una o dos cuerdas no contiene el soporte
de ninguna dependencia. Añadir \(\min(\beta_g,2)\) cuerdas da el rango indicado. Restándolo
de \(2v\) y sumando sobre componentes sale la última desigualdad.
\end{proof}

\begin{lema}[Pesos y número de tipos]\label{lem:tipos-llegada}
Póngase \(L_s=v+c_s-D\) y \(d_i=\max_{f\ni i}|f\cap U|\). Para \(v,c_s,L_s\) fijos,
la suma sobre tipos etiquetados realizables satisface
\[
 \frac1{v!}\sum_{\rm tipos}\prod_i d_i\le \frac{C^v(4v^2)^{L_s}}{c_s!}.
\]
Supóngase que los únicos conflictos diagonales del soporte son dos pares en líneas distintas.
Con cuatro sitios, las cotas \((D,p)\) para cero a cuatro coincidencias axiales son
\((6,7),(5,6),(4,5),(3,4),(3,4)\), cuando la clase es realizable.
Con tres sitios, los pares comparten un sitio y las cotas son \((4,5),(3,4)\)
según sus extremos tengan cero o una coincidencia axial.
\end{lema}
\begin{proof}
Para cada sitio se elige la primera fibra en que alcanza su máximo \(d_i\), en un orden
fijo de familias. Unimos los \(a_f>0\) sitios asignados a cada fibra \(f\) con un nodo de
su color. Añadimos aristas de igualdad entre sitios hasta conectar cada componente sin
formar ciclos. Las estrellas admiten un árbol equivalente de igualdades entre sitios.
Este bosque tiene rango \(v-c_s\); para completar la base hacen falta \(L_s\) aristas,
cada una con a lo sumo \(4v^2\) elecciones. La base recupera las cuatro particiones:
una igualdad en su espacio generado se cumple en la realización y por ello pertenece al
tipo. Enraizamos en un sitio por componente y reunimos las raíces bajo un nodo común.
Hay a lo sumo \(2v+1\) nodos, un número fijo de colores y \(v\) etiquetas de sitios.
La codificación por árboles planos tiene a lo sumo \(C^v v!\) posibilidades, por la cota
de Catalan. Cada tipo admite al menos \(c_s!\prod_f(a_f-1)!\) órdenes de hijos distintos.
Los miembros de cada estrella se reconocen por sus etiquetas. Con \(D_f=|f\cap U|\),
\[
 \prod_f\frac{D_f^{a_f}}{(a_f-1)!}\le 2^v\exp\!\left(\sum_fD_f\right)
 \le 2^ve^{4v}.
\]
Así se obtiene la suma ponderada. Al contar los soportes sin orden, cada soporte da \(v!\)
etiquetados, incluso cuando varios de ellos tienen el mismo tipo.

En la clasificación de los dos conflictos se incluyen todas las coincidencias axiales.
Por hipótesis los dos pares agotan las aristas diagonales. Sin coincidencias axiales, los
cuatro sitios forman dos árboles. La primera
coincidencia los conecta, la segunda crea un ciclo, y desde la tercera la cota general da
\(p=4\). El lema de rango reduce la dimensión a \(6,5,4\) en los tres primeros casos.
Con tres coincidencias están fijadas las dos filas y una columna del segundo par dadas las del
primero, o viceversa. Quedan a lo sumo tres parámetros, aun si hay una cuarta coincidencia.
En la cuña de tres sitios el primer grafo es un árbol y la coincidencia entre extremos crea
un ciclo. Si hay un tercer conflicto diagonal se vuelve a contar el número de ciclos en el
grafo completo. El lema de rango y la Proposición~\ref{prop:soportes-llegada} dan entonces
\(D\le p-1\) con el exponente de ese grafo. Los pares en el mismo
cuadrado se eliminan por la sustracción de su error.
\end{proof}
